\chapter{Coronary Artery Bypass Graft}
\index{Coronary Artery Bypass Graft}

\section*{Definition and Overview}
Coronary artery bypass grafting (CABG), commonly called heart bypass surgery, is an operation that creates new routes for blood to flow around narrowed or blocked coronary arteries. Healthy blood vessels taken from elsewhere in the body -- most often a leg vein, an internal chest artery, or a forearm artery -- are used to bypass the blockages, restoring blood supply to the heart muscle. Bypass surgery is a major operation performed under general anaesthetic with the heart either stopped (using a heart-lung machine) or beating. It is one of the most effective treatments for extensive coronary artery disease and can relieve angina, improve quality of life, and extend life in selected patients.

\section*{Epidemiology}
Coronary artery disease is the single largest cause of death worldwide, and bypass surgery is one of the most commonly performed major operations. In many countries, thousands of bypass procedures are performed each year. The typical candidate has disease in several coronary arteries, narrowing of the left main artery, or diabetes with extensive disease, where surgery has been shown to offer better outcomes than stenting. Outcomes are generally excellent: operative mortality is low in fit patients, and grafts remain functional for many years, particularly the internal chest artery graft.

\section*{Aetiology and Pathophysiology}
Bypass surgery treats the consequences of atherosclerosis, the buildup of fatty plaque in the coronary artery walls.

\begin{itemize}
    \item \textbf{Atherosclerosis:} cholesterol-rich plaques narrow and stiffen the coronary arteries over decades
    \item \textbf{Plaque rupture:} a torn plaque triggers a blood clot, causing heart attack
    \item \textbf{Stable angina:} fixed narrowings limit blood flow when the heart works harder
    \item \textbf{Left main disease:} narrowing of the artery supplying most of the left heart, carrying high risk
    \item \textbf{Three-vessel disease:} extensive disease in the three main coronary arteries
    \item \textbf{Diabetes:} accelerates atherosclerosis and often produces more diffuse disease
    \item \textbf{Prior stents:} grafts may be needed when stents are unsuitable or have failed
\end{itemize}

\section*{Clinical Features}
The symptoms that lead to bypass surgery reflect reduced blood flow to the heart.

\begin{itemize}
    \item Chest pain, pressure, or heaviness with exertion (angina)
    \item Shortness of breath with activity
    \item Fatigue and reduced exercise capacity
    \item A prior heart attack
    \item An abnormal stress test or coronary angiogram showing severe disease
    \item Symptoms poorly controlled despite medical therapy
\end{itemize}

\section*{Differential Diagnosis}
The decision for surgery is based on symptoms, the results of angiography, and the extent of disease, and surgery is compared against stenting and medical therapy.

\begin{itemize}
    \item Percutaneous coronary intervention (stenting) as an alternative revascularisation strategy
    \item Medical therapy alone, in people with mild or limited disease
    \item Non-cardiac causes of chest pain, such as reflux or musculoskeletal pain
    \item Heart failure from other causes, such as cardiomyopathy or valve disease
    \item Aortic stenosis or other valve disease needing its own treatment
\end{itemize}

\section*{When to Seek Medical Care}
If you have angina or have been offered bypass surgery, you will be assessed by a cardiologist and a cardiac surgeon. Urgent treatment is needed for unstable symptoms.

\textbf{Contact your local emergency services for:}
\begin{itemize}
    \item Chest pain at rest, or pain lasting more than a few minutes
    \item Pain spreading to the arm, jaw, neck, or back
    \item Shortness of breath, sweating, nausea, or dizziness with chest pain
    \item Signs of a heart attack, or symptoms that recur after surgery
\end{itemize}

\section*{Diagnosis and Investigations}
Assessment before surgery confirms the diagnosis and ensures the patient is as fit as possible.

\begin{itemize}
    \item \textbf{Coronary angiography:} the essential test defining which arteries are blocked
    \item \textbf{Echocardiography:} assessment of heart pump function (ejection fraction) and valves
    \item \textbf{ECG and Holter monitoring:} recording of heart rhythm and evidence of prior heart attack
    \item \textbf{Blood tests:} kidney function, liver function, full blood count, and diabetes checks
    \item \textbf{Carotid and peripheral artery assessment:} checking for other vascular disease
    \item \textbf{CT of the chest:} to assess the aorta and plan the operation
    \item \textbf{General health optimisation:} smoking cessation, blood pressure and diabetes control, and dental review
\end{itemize}

\section*{Management}
Bypass surgery is performed by a specialised cardiac surgical team.

\begin{itemize}
    \item \textbf{Pre-operative preparation:} medicines adjusted, smoking stopped, and fitness optimised
    \item \textbf{Anaesthesia:} general anaesthetic with careful monitoring
    \item \textbf{Harvesting grafts:} the internal chest (mammary) artery, leg vein, or forearm artery is prepared
    \item \textbf{On-pump surgery:} the heart is stopped and blood is circulated by a heart-lung machine
    \item \textbf{Off-pump surgery:} the operation is performed on the beating heart in selected patients
    \item \textbf{Grafting:} grafts are attached beyond each significant blockage
    \item \textbf{Hospital stay:} typically five to seven days, with time in intensive care first
    \item \textbf{Cardiac rehabilitation:} supervised exercise and education over the following weeks
    \item \textbf{Long-term medicines:} antiplatelet agents, statins, and blood pressure control
\end{itemize}

\section*{Complications}
Bypass surgery is a major operation, and although outcomes are generally excellent, complications can occur.

\begin{itemize}
    \item Bleeding requiring re-operation
    \item Infection of the chest wound or leg wound
    \item Irregular heart rhythms, especially atrial fibrillation
    \item Stroke or transient neurological problems
    \item Kidney injury
    \item Heart attack or heart failure around the time of surgery
    \item Lung complications such as pneumonia or fluid collections
    \item Graft failure in the months or years after surgery
    \item Post-operative cognitive changes, usually temporary
    \item Memory or concentration problems in some patients
\end{itemize}

\section*{Prognosis and Outcomes}
For suitable patients, bypass surgery reliably relieves angina and improves exercise capacity, and in those with extensive disease it improves long-term survival compared with medical therapy or stenting. The internal chest artery graft has a very high ten-year patency rate, while vein grafts remain open for a shorter period, which is why risk factor control is essential. Most people return to normal activities within six to twelve weeks. Recurrent angina can develop years later as grafts narrow or native disease progresses, but many patients remain symptom-free for a decade or more.

\section*{Prevention and Self-Care}
\begin{itemize}
    \item Stop smoking before surgery and remain a non-smoker afterwards
    \item Take antiplatelet medicines, statins, and blood pressure medicines as prescribed
    \item Attend cardiac rehabilitation and build up activity gradually
    \item Follow a heart-healthy diet and keep cholesterol and blood sugar controlled
    \item Achieve and maintain a healthy weight
    \item Attend regular follow-up with the cardiologist and GP
    \item Protect the sternum: avoid heavy lifting and straining during recovery
    \item Report new chest pain, fever, or wound problems promptly
\end{itemize}

\section*{Sources and Further Reading}
The following reputable sources provide further information for healthcare students and patients seeking reliable, current advice about coronary artery bypass surgery.

\begin{itemize}
    \item NHS. \textit{Coronary artery bypass graft (CABG).} https://www.nhs.uk/conditions/coronary-artery-bypass-graft-cabg/
    \item Heart Foundation Australia. \textit{Bypass surgery.} https://www.heartfoundation.org.au/
    \item Mayo Clinic. \textit{Coronary bypass surgery: overview and recovery.} https://www.mayoclinic.org/tests-procedures/coronary-bypass-surgery/
    \item MedlinePlus. \textit{Coronary artery bypass surgery.} https://medlineplus.gov/ency/article/002946.htm
    \item MSD Manual. \textit{Coronary artery bypass grafting: professional reference.} https://www.msdmanuals.com/
    \item American Heart Association. \textit{Bypass surgery (CABG).} https://www.heart.org/
\end{itemize}
